\section{Design}
\label{sec:design}

We designed \sysname, a serving layer for structured LLM agents that exploits the heterogeneity of compiler-generated requests to schedule the KV cache more efficiently, closing the gap between the runtime that generates prompts and the engine that serves them. When a developer declares an LLM call in code, the framework's compiler or runtime lowers the call into a prompt whose layout it fixes: a header carrying the signature, docstrings, type schemas and output format, followed by one region per bound argument. The runtime knows which of these regions is a program constant, which is carried over from an earlier call, and which is fresh. The engine sees only bytes. \sysname recovers that structure on the engine side, learns from it how each component of a request lives and dies, and schedules the KV cache by that knowledge instead of by recency alone.

\sysname has two planes (Figure~\ref{fig:overview}). The \emph{discovery plane} lifts each incoming prompt back to its call-site signature and bindings, then combines this metadata with request history to learn how each binding evolves and how the agent moves between call sites. Its output is an intermediate representation of every request, $\langle$call site, bindings, provenance, lifetime, sharing, future reuse$\rangle$, together with a time-annotated probabilistic finite-state machine (TPFM) of the application. Both describe properties the program fixes rather than statistics of a particular input, so they converge after a handful of observations per call site and are then applied exactly. The \emph{scheduling plane} consumes this knowledge at two levels. Per request, it re-lays out the prompt so that stable and shared components lead and volatile ones trail, which enlarges the prefix that consecutive calls share. Across requests, it predicts the prompts a session will send next, promotes their prefixes from host to device before they are needed, and orders eviction by lifetime and predicted reuse, so that a reusable KV segment resides in the tier where it will be reused.

\subsection{Discovering Heterogeneity}
\label{sec:discovery}

\sysname first lifts each received prompt back to its call-site metadata. It then analyzes how each binding evolves across the request history along two dimensions, intra-call-site evolution and cross-call-site sharing, and assigns every binding a lifetime and sharing class. Separately, it learns the application's call structure as a TPFM. Lifting happens on every request; the analysis runs off the request path, when a session completes. The scheduling plane only reads the result.

\subsubsection{Lifting prompts back into metadata}
\label{sec:lifting}

Even though the lowering is invisible to the engine once a prompt crosses the boundary between runtime and serving engine, the generated request still retains recoverable structure, because a framework lowers every call of the same call site in the same way. Two properties of the lowering are all that lifting relies on. First, the static and the dynamic parts of the prompt are separable: the header, schema rows and output-format tail are emitted from the declaration and never change between calls of a site, while every bound argument is rendered as its own contiguous block, labelled by the parameter name. Second, the order and labelling of those blocks follow the declaration, so a block can be attributed to a parameter without understanding its content. Figure~\ref{fig:heterogeneity} shows this shape for the \texttt{implement(spec, plan, advice, attempts)} site of a coding agent, and the same shape appears when the same call is declared as a DSPy signature: a framework-static system prompt, one labelled field per argument in declaration order, and a framework-static tail. Lifting is therefore a per-framework grammar rather than a learned model, and its value lies in the abstraction level of its output: it turns a stream of opaque requests into a stream of \emph{(call site, bindings)} pairs, which is what lets the rest of the design attribute heterogeneity to program structure rather than to string statistics.

Algorithm~\ref{alg:decompile} gives the lifting grammar. It strips the fixed output-format tail, locates the header (which \sysname itself may have moved behind the values, Section~\ref{sec:relayout}), reads the parameter names in declaration order from the signature, and scans the value region for binding blocks, gluing continuation lines of a multi-line value to the binding they belong to. A call site is identified by the static part only: the signature, the parameter names with their declared types and descriptions, and the model it targets. Everything that varies between calls, the bindings and later the reply, is attached to the observation of that call.

One logical call may span several requests: a tool-using call re-sends the transcript after every tool result, and a typed-output failure re-sends the call with feedback. \sysname folds these into one observation by checking whether the new transcript extends the previous one, so the TPFM sees logical calls, and each observation records its number of turns and the gaps between them.

\subsubsection{Identifying heterogeneous evolution patterns}
\label{sec:evolution}

Lifting says what a request is made of; the traces say how each part behaves. For every binding of every call site, \sysname maintains three counters, updated once per completed session by walking its calls in order (Algorithm~\ref{alg:heterogeneity}).

\paragraph{Provenance.} For each binding value, \sysname asks which earlier value of the same session explains it. The same name with an equal value at an earlier call is a \emph{copy}. A value that starts with the latest earlier value of that name minus its closing delimiter is an \emph{extension}: a list, dict or string that grew by appending. A value that appears inside an earlier reply, under another binding name, or as an option an earlier call selected, is \emph{derived} from that call. A value none of these explain is recorded as the literal itself. Because that counter is keyed by the literal, a literal accumulates support only when the same bytes recur across sessions, which is what distinguishes a program constant from a fresh input.

\paragraph{Intra-call-site evolution.} When a session revisits a call site, each binding is compared with its value at the latest earlier visit and the transition is counted as \emph{same}, \emph{extend} or \emph{fresh}. This is the evidence that separates a value fixed for the session from one that accumulates, and both from one that is replaced on every visit. In the example of Figure~\ref{fig:heterogeneity}, \texttt{spec} is the same on every visit, \texttt{attempts} extends, and \texttt{advice} is fresh.

\paragraph{Cross-call-site sharing.} A binding is marked \emph{cross} when another call site earlier in the same session already carried the same value, or a prefix that the current value extends; otherwise it is \emph{own}. Sharing decides whether two sibling call sites of one session can be made to share a radix path.

A provenance rule becomes dominant once it has at least two observations and accounts for at least half of all observations of the binding. From the dominant rule and the transition counts, every binding is classified as one of four evolution patterns. \emph{Constant}: the same literal recurs across sessions. \emph{Copied}: the value is reused from an earlier call in the same session and holds between visits. \emph{Extended}: each new value grows from the previous one by appending. \emph{Volatile}: the value is fresh on most calls and is not reused. Concretely, a constant ranks first; otherwise a binding is scored by the fraction of visits on which it was fresh, tie-broken by its most common non-fresh transition (same before extend); a binding a session visits only once has no transition evidence and falls back to the rank of its provenance (copy before extend before anything else). Derived values rank as volatile for layout, since they change with every reply, but they remain reconstructible from the reply that produced them, which Section~\ref{sec:anticipate} exploits. Orthogonally, a binding is \emph{cross-stable} when at least half of its observations are cross and it is not volatile: it either has no transition evidence or is fresh on at most half of its visits.

The classification of a site is frozen once every one of its bindings has been observed at least eight times, and never changes afterwards: each change of order breaks the prefix once, so the order is decided once. The frozen layout is a stable sort of the declared order by (cross-stable first, then the stability score), which keeps the declared order among ties. Alongside it, \sysname records how the message is assembled (binding order and the header text for each value shape the framework renders differently) and counts how often the learned rules rebuild a real message byte for byte. That count gates speculation: a reconstructed prompt is trusted only for a site whose rules have reproduced a real message at least once. With this the per-request record is complete: call site and bindings from lifting, provenance, lifetime and sharing from the counters, and future reuse from the TPFM described next.

\subsubsection{Deriving the application-level agent TPFM}
\label{sec:tpfm}

The classes say what will be needed again; the TPFM says when. \sysname models an agent as $\mathcal{A} = (Q, P, q_0, F, \mathcal{T})$, where each state $q \in Q$ is a call site, $P$ is the transition probability between call sites, $q_0$ is the entrance call site, $F(q)$ is the probability that the agent terminates at $q$ with $F(q) + \sum_{q'} P(q, q') = 1$, and $\mathcal{T} = \{D^{\mathrm{exec}}_q, D^{\mathrm{gap}}_{q,q'}\}$ annotates each state with the execution-time distribution of its call and each transition with the distribution of the gap between the completion of $q$ and the invocation of $q'$.

A program is identified by its entrance call site: a session is bound to the program whose entrance matches its first call, and a session that joins mid-program is served but excluded from learning. From the engine's side a session is only a chain of lifted calls, and each completed chain contributes its sequence of call sites, the gap after each call, and each call's engine time, turn count and tool gaps. Chains are recorded in a prefix tree with branching frequencies and timing samples on its edges. The tree is folded into the state machine by an Alergia-style merge, run when the program has seen eight sessions and at every doubling thereafter: two states with the same call site are merged when their termination rates and every successor rate are statistically indistinguishable under a Hoeffding bound with $\alpha = 0.05$. After a fold, every recorded session is replayed onto the merged states so that state-specific gap and duration samples survive.

Transition frequencies of a freshly merged machine are unreliable while few sessions have been seen, so $P$ is estimated by a bounded-context predictor rather than read off the machine: a trie over call-site subsequences of length at most four, smoothed by Witten--Bell backoff across the suffixes of the current history. The machine supplies $\mathcal{T}$: the gap distribution on the edge the session is about to take and the duration distribution of the state it will enter, with per-pair and per-site statistics as fallback when a state has no samples. Prediction is a bounded fan-out search over this pair. From the session's history it expands the three most likely continuations at each step, prunes paths below probability $0.02$ or beyond a 120-second horizon, and stops after 64 expansions or eight steps. Each predicted call site receives the probability mass of every path that reaches it, an earliest median arrival $t_{50}$, and a pessimistic arrival $t_{90}$ obtained by accumulating the spread of the gap and duration quantiles along the path. The same predictor yields $F$ for the current history, which the server uses to shorten the idle timeout of a session that has probably finished.

\subsection{Heterogeneity-aware KV Cache Scheduling}
\label{sec:scheduling}

Scheduling the KV cache involves two layers of management. Prompt re-layout acts on each request as it passes through: it reorders the bindings of a call site by lifetime and moves shared components to the front, so that the prefix a later call can reuse is as long as the data allows. KV management acts across requests: it predicts the next call's prompt from the learned rules and the TPFM, prefetches its host-resident prefix by deadline, protects prefixes with predicted reuse, and evicts what has no further use.

On the request path, \sysname lifts each arriving request, binds its session to a program and, if the site's classification is frozen, re-emits the call message in the frozen order before rendering. A continuation turn of an open call is re-emitted in the same order and enters the engine with a higher scheduling priority than a fresh call, so that a tool loop in flight is not starved. A fresh call advances the session: its epoch is bumped, which voids every job planned for the previous step, and its successors are planned immediately with the arrival time as anchor, so that their preparation overlaps this call's own execution. When the reply returns, the planner is told which part of the served prompt is one-off, the anchor moves to the true completion time, and the plan is refreshed. A session ends on an explicit close or an idle timeout; its private cache is then retired and its observations are folded into the program.

\subsubsection{Prompt re-layout: stable first, shared first}
\label{sec:relayout}

KV reuse requires an identical token prefix. In the order the program declares its parameters, a volatile binding may precede a stable one; the shared prefix of two consecutive calls then ends at the first volatile byte, and everything behind it, however stable, is recomputed. In Figure~\ref{fig:heterogeneity}, \texttt{plan} and \texttt{advice} precede \texttt{attempts}, so whenever they change the reusable prefix of \texttt{attempts} is contaminated. \sysname re-emits the call message in the frozen order instead: constant, then copied, then extended, then volatile. Binding blocks move as units, every block and the fixed tail are preserved byte for byte, the transformation is idempotent, and a message whose blocks are not contiguous is left untouched. The longest common prefix between calls of one site now ends at the end of the accumulating history, which by construction is a prefix of its own next value.

If a binding is shared by several call sites (cross-stable), it moves to the front and the static header behind it, so that different call sites of one session share one prefix and only their trailing headers differ. This is what makes sharing a first-class dimension: sibling call sites that carry the same session state, such as a claim and its accumulated evidence, are served from one radix path. The decision is made per site at freeze time, from whether the site's leading binding is cross-stable. Re-layout changes the order in which the model reads components, not their content; its effect on task quality is measured in Section~\ref{sec:eval}.

\subsubsection{Anticipating the next prompt}
\label{sec:anticipate}

Re-layout makes the reusable prefix long; KV management makes sure it is on the device when the call arrives. For each call the TPFM predicts, \sysname rebuilds the prompt that call will send from what the live session already carries, applying each binding's dominant provenance rule in the frozen order: a copy takes the latest value of that name, a constant takes its literal, a derived value is read back from the reply or binding it came from, and an extended value contributes the known head of its next value, after which reconstruction stops because the tail is not yet known. The rebuilt text is rendered through the same chat template as a real request and tokenized. If the whole message resolved and the site's rules have been verified byte-exact before, the target is the complete prompt; otherwise it is the rendered prefix up to the end of the reconstructed head, or, failing that, the site's static head alone. Targets shorter than a cache block are dropped.

Every target is priced against a shadow ledger of the two-tier cache. The ledger mirrors the radix cache's content addressing in stride-sized chain-hashed blocks, keeps a device tier in LRU order and an inclusive host tier, and is recalibrated by every real request's measured device and host hits, so its drift is bounded by one request. It splits a target into tokens already on the device, host-resident tokens a promotion would cover, and tokens that must be computed. A target becomes a \emph{promote} job when its uncomputed part is within one stride and it has host-resident tokens, and a \emph{hold} job otherwise. A hold job only carries eviction protection, but it is upgraded to a promotion as soon as the ledger learns that more of its prefix is cached, for instance because another session computed the shared head. Each job carries the call's probability, a value equal to that probability times the exposed prefill it saves, and two times derived from the TPFM: a deadline $\max(t_{\mathrm{now}},\, a + t_{50} - \tfrac{1}{2}(t_{90}-t_{50}))$, where $a$ is the session's anchor, and a pessimistic arrival $\max(d,\, a + t_{90})$ after which the job is stale and discarded.

\subsubsection{Prefetch by deadline}
\label{sec:deadline}

The planner loop wakes on every plan change and at least every 15 ms. It discards stale and voided jobs, pushes the current eviction map to the engine, then executes at most one promotion. A job is released when the engine is idle, or when the time approaches its latest start: the deadline minus the profiled cost of computing its uncached tokens and promoting its host tokens at the device's measured rates. Among released jobs whose cached frontier reaches past what is already on the device, the one with the earliest deadline runs, value breaking ties. A promotion is an RPC into the engine's scheduler, not a request. The modified hierarchical cache starts the host-to-device copy on its background promotion thread and lowest-priority stream, locks the loaded nodes until the copy lands, and lets a batch that later admits them wait on the copy per layer. A promotion may use free device slots and reclaim retired cache but never evicts cache a live session might reuse, and it keeps a prefill chunk of space in reserve so that concurrent promotions cannot lock the pool against real prefill. A promotion refused for space pauses all promotions for a few cycles; one deferred behind real work is retried; a job whose RPC fails three times is dropped.

\subsubsection{Protect reuse, evict volatile}
\label{sec:eviction}

The device tier is evicted in four bands, lowest first and LRU inside a band: retired $<$ transient $<$ unclassified $<$ planned. The bands follow the lifetime of the bytes they hold. \emph{Retired} holds the private cache of sessions that have ended; no live session can reuse it, so it is reclaimed first. \emph{Transient} holds the one-off tail of each served prompt, everything past the head the provenance rules can rebuild: the volatile bindings and the generated tokens. The reconstructible head stays in the \emph{unclassified} band with all cache no plan covers. \emph{Planned} prefixes are ranked by predicted reuse time: a prefix scores $p \cdot \exp(-\max(0, d - t)/\tau)$ with $\tau = 60$~s, so a prefix needed sooner is protected more than an equally likely prefix needed later, and a radix node covered by several sessions' plans receives the sum of their scores, so a head shared across sessions outranks any single session's private history. The map is rebuilt from scratch on every push and applied after the previous protection has been undone, so no protection outlives the plan that granted it. Demotion to transient happens when a call is served; retirement when its session ends.

The host tier stays LRU as the safety net: it evicts retired cache first and everything else by recency, and the plan never reorders it. Applying the prediction policy to both tiers would let one misprediction remove a prefix from the device and the host at once, turning a recoverable device miss into a full recompute. Leaving the host recency-ordered keeps every prefix the device gave up reloadable until it ages out.
